\documentclass[12pt]{article}

\usepackage{bbm}
\usepackage{fancyhdr}
\usepackage{titling}
\usepackage{enumitem}
\usepackage{amsmath}
\usepackage{dsfont}
\usepackage{amsfonts}
\usepackage{amssymb}
\usepackage{amsthm}
\usepackage{geometry}
\usepackage{tikz}
\usepackage{algorithm}
\usepackage{algpseudocode}
\geometry{verbose,tmargin=1in,bmargin=1in,lmargin=1in,rmargin=1in}
\renewcommand{\vec}[1]{\mathbf{#1}}
\newcommand{\dec}{\text{dec}}
\newcommand{\divi}{\mathbin{\textnormal{div}}}
\title{Homework 1}
\author{Satvik Prasad}
\date{}
\newtheorem{lemma}{Lemma}[section]
\fancypagestyle{plain}{%
  \fancyhf{}
  \fancyhead[L]{\theauthor}
  \fancyhead[R]{\thedate}
}
\pagestyle{plain}
\makeatletter
\def\@maketitle{%
  \newpage
  \null
  \vskip 1em%
  \begin{center}%
    \let \footnote \thanks
    {\LARGE \@title \par}%
    \vskip 1em%
  \end{center}%
  \par
\vskip 1em}
\makeatother
\usepackage{cmbright}

\newcommand{\mat}[1]{\mathbf{#1}}

% ---------- Probability theory ----------
% Core operators (optional argument-free; use as \E{X}, \Var{X}, \Cov{X}{Y})
\DeclareMathOperator{\EOp}{\mathbb{E}}
\DeclareMathOperator{\VarOp}{Var}
\DeclareMathOperator{\CovOp}{Cov}
\DeclareMathOperator{\CorrOp}{Corr}

\renewcommand{\P}[1]{\mathbb{P}\left(#1\right)}        % \P{A}
\newcommand{\E}[1]{\EOp\left[#1\right]}                % \E{X}
\newcommand{\Var}[1]{\VarOp\left(#1\right)}            % \Var{X}
\newcommand{\Cov}[2]{\CovOp\left(#1,\,#2\right)}       % \Cov{X}{Y}
\newcommand{\Corr}[2]{\CorrOp\left(#1,\,#2\right)}     % \Corr{X}{Y}
\newcommand{\SD}[1]{\operatorname{SD}\left(#1\right)}  % \SD{X}

% Subscripted versions, e.g. \Esub{\theta}{X}, \Psub{\theta}{A}
\newcommand{\Esub}[2]{\EOp_{#1}\left[#2\right]}
\newcommand{\Psub}[2]{\mathbb{P}_{#1}\left(#2\right)}

% Conditioning, indicators, independence
\newcommand{\given}{\,\middle|\,}                      % \E{X \given Y}
\newcommand{\ind}[1]{\mathbbm{1}\left\{#1\right\}}     % \ind{X > 0}
\newcommand{\indep}{\perp\!\!\!\perp}                  % X \indep Y
\newcommand{\iid}{\overset{\text{iid}}{\sim}}          % X_i \iid \Normal(0,1)

% Convergence
\newcommand{\convp}{\xrightarrow{\;p\;}}               % in probability
\newcommand{\convd}{\xrightarrow{\;d\;}}               % in distribution
\newcommand{\convas}{\xrightarrow{\text{a.s.}}}        % almost surely
\newcommand{\convLp}[1]{\xrightarrow{L^{#1}}}          % in L^p

% Probability spaces
\newcommand{\F}{\mathcal{F}}                           % sigma-algebra
\newcommand{\B}{\mathcal{B}}                           % Borel sigma-algebra
\newcommand{\pspace}{(\Omega, \F, \mathbb{P})}

% Common distributions
\newcommand{\Normal}{\mathcal{N}}
\newcommand{\Bern}{\operatorname{Bernoulli}}
\newcommand{\Bin}{\operatorname{Binomial}}
\newcommand{\Pois}{\operatorname{Poisson}}
\newcommand{\Unif}{\operatorname{Uniform}}
\newcommand{\Exp}{\operatorname{Exp}}
\newcommand{\Geom}{\operatorname{Geometric}}
\newcommand{\Gam}{\operatorname{Gamma}}
\newcommand{\Beta}{\operatorname{Beta}}

% Number sets
\newcommand{\R}{\mathbb{R}}
\newcommand{\N}{\mathbb{N}}
\newcommand{\Z}{\mathbb{Z}}


\begin{document}
\begin{enumerate}
    \item We know $X, Y$ are independent random variables such that $X, Y \sim \text{Exp}(\lambda).$ We want to find the joint density of $U = X, V = X + Y.$

    Define the map $T: \mathbb [0, \infty)^2 \to D$ such that
    $$
    T(x, y) = (x, x + y),
    $$
    with inverse $T^{-1}: D \to [0, \infty)^2$ $$
    T^{-1}(u, v) = (u, v - u).
    $$

    For our domain, we need $u \geq 0$ and $v \geq u,$ meaning $D = \{(u, v) : u \geq 0, v \geq u\}$ as required.

    This map has Jacobian $$
    J = \begin{vmatrix}
        1 & 0 \\
        -1 & 1
    \end{vmatrix} =  1,
    $$
    and thus our joint density for $U, V$ satisfies
    $$
    f_{U, V}(u, v) = f_{X, Y}(u, v - u).
    $$
    the joint density for $X, Y$ is the following, since $X, Y$ independent:
    $$
    f_{X, Y}(x, y) = \begin{cases}
        \lambda^2e^{-\lambda( x + y)} & x \geq 0, y \geq 0 \\
        0 & \text{otherwise}
    \end{cases},
    $$
    meaning we get 
    $$
    f_{U, V}(u, v) = \begin{cases}
        \lambda^2 e^{-\lambda v} & v \geq u \geq 0,  \\
        0 & \text{otherwise}
    \end{cases}.
    $$

    Then, we want to find the conditional density of $U \mid V = a,$ which satisfies
    $$
    f_{U \mid V = a}(u) = \dfrac{f_{U, V}(u, a)}{f_V(a)}.
    $$
    Computing the marginals, when $v \geq 0$ we have
    $$
    f_V(v) = \int_{0}^{v} \lambda^2 e^{-\lambda v} \, du = v\lambda^2 e^{-\lambda v},
    $$
    yielding conditional density $$
    f_{U \mid V = a}(u)  = \dfrac{\lambda^2 e^{-\lambda a}}{a \lambda^2 e^{-\lambda a}} = \dfrac{1}{a}
    $$
    when $0 \leq u \leq a,$ and $0$ otherwise. Therefore, $(U \mid V = a) \sim \text{Uniform}(0, a)$ as required.

    \item Computing the MGF of $X \sim \text{Uniform}(a, b),$ we have
    $$
        M(s) = \int_{a}^{b} \dfrac{e^{st}}{b - a} \, dt = \dfrac{e^{bs} - e^{as}}{s(b- a)} = \dfrac{1}{b-a} (s^{-1}e^{bs} - s^{-1} e^{as}).
    $$

    To find the $k$-th moment, we can write $$
    e^{bs} - e^{as} = \sum_{j = 1}^{\infty} \dfrac{b^j - a^j}{j!} s^j, 
    $$
    yielding
    $$
    M(s) = \dfrac{1}{b-a} \sum_{j = 0}^{\infty} \dfrac{b^{j + 1} -a^{j + 1}}{(j + 1)!}s^{j}.
    $$
    This yields $k$-th derivative
    \begin{align*}
    M^{(k)}(s) &= \dfrac{1}{b-a} \sum_{j = k}^{\infty} \dfrac{b^{j + 1} - a^{j + 1}}{(j+1)!} \dfrac{j!}{(j - k)!} s^{j - k} \\
    &= \dfrac{1}{b-a} \sum_{j = 0}^{\infty} \dfrac{b^{j +k +  1} - a^{j +k +  1}}{(j+k  + 1)(j)!} s^{j},
    \end{align*}
    meaning $M^{(k)}(0)$ satisfies
    $$
    M^{(k)}(0) =  \dfrac{b^{k+1} - a^{k+1}}{(b-a)(k + 1)} = \mathbb E[X^k]
    $$
    as required.

    \item Define $S = X_1 + \ldots + X_n$ where $N$ is an RV with PGF $G(s).$ Conditioning on $N,$ we have
    $$
    \mathbb P(S \leq s \mid N  = n) = \mathbb P(X_1 + X_2 + \ldots + X_n \leq s).
    $$
    Thus, we have $$
    \mathbb P(S \leq s) = \sum_{n = 0}^{\infty} \mathbb P(X_1 + \ldots + X_n \leq s) \mathbb P(N = n) = \sum_{n = 0}^{\infty} \mathbb P(S_n \leq s) \mathbb P(N = n).
    $$
    Consider the random variable $S_n = X_1 + \ldots + X_n.$ The MGF of $S_n$ satisfies $M_n(s) = \mathbb E[e^{sS_n}] = M(s)^{n}.$

    By the above, we know
    $$
    f_S(s) = \sum_{n = 0}^{\infty} f_{S_n}(s) \mathbb P(N = n).
    $$
    Computing our MGF, we have
    \begin{align*}
        M_S(t) &= \int_{-\infty}^{\infty} e^{st} \sum_{n = 0}^{\infty} f_{S_n}(s) \mathbb P(N = n) \, ds \\
        &= \sum_{n = 0}^{\infty} \mathbb P(N = n)\int_{\mathbb R} e^{st} f_{S_n}(s) \, ds \\
        &= \sum_{n = 0}^{\infty} \mathbb P(N = n) M_{S_n}(t) \\
        &= \sum_{n = 0}^{\infty} \mathbb P(N = n) M(t)^n \\
        &= G(M(t))
    \end{align*}
    as required.

    \item \begin{enumerate}
        \item  We know $\Cov{X}{Y} = \mathbb E[(X - \mu)(Y - \lambda)]$ where $\mu = \mathbb E[X]$ and $\lambda = \mathbb E[Y].$ By the linearity of expectation, 
        $$
        \mathbb E[X] = \sum_{k = 1}^{n} \mathbb E[X_k], \qquad \mathbb E[Y] = \sum_{k = 1}^{n} \mathbb E[Y_k].
        $$

        Substituting, we get
        $$
        \Cov{X}{Y} = \mathbb E[XY] - \mathbb E[X] \mathbb E[Y] = \mathbb E\left[\sum_{k = 1}^{n}\sum_{j = 1}^{n} X_i Y_j\right] - \sum_{k = 1}^{n} \sum_{j = 1}^{n} \mathbb E[X_k] \mathbb E[Y_j].
        $$
        Since $X_i$ and $Y_j$ are independent when $i \neq j,$ we can write 
        \begin{align*}
        \Cov{X}{Y} &= \sum_{k = 1}^{n} \sum_{j = 1}^{n} \mathbb E[X_i Y_j] - \mathbb E[X_k] \mathbb E[Y_j] \\
        &= \sum_{k = 1}^{n} \mathbb E[X_k Y_k] - \mathbb E[X_k] \mathbb E[Y_k] = \sum_{k = 1}^{n}\Cov{X_k}{Y_k}
        \end{align*}
        as required.

        \item Let $X_k$ be the number of points scored by player $A$ in game $k,$ and equivalently for $Y_k.$ Then, $X_i$ and $Y_j$ are independent whenever $i \neq j,$ meaning if $A = X_1 + \ldots + X_n$ and $B = Y_1 + \ldots + Y_n,$ we have
        \begin{align*}
        \Cov{A}{B} &= \sum_{i = 1}^{n} \Cov{X_i}{Y_i} \\
        &= \sum_{i = 1}^{n} 2pq - (2p^2 + 2pq)(2q^2 + 2pq) \\
        &= \sum_{i = 1}^{n} 2pq - 4pq = -2npq
        \end{align*}
        as required.
    \end{enumerate}

    \item \begin{enumerate}
    
    \item Partition on $X > Y,$ yielding \begin{align*}
    \mathbb P(\min(X, Y) + \max(X, Y) \leq k) &= \mathbb P( \min(X, Y) + \max(X, Y) \leq k \cap {X > Y})  \\
    &+ \mathbb P( \min(X, Y) + \max(X, Y) \leq k \cap {X \leq Y}) \\
    &= \mathbb P(Y + X\leq k \cap X > Y) + \mathbb P(X + Y \leq k \cap X \leq Y).
    \end{align*}
    Since these events are disjoint, we therefore know
    \begin{align*}
    \mathbb P(\min(X, Y) + \max(X, Y) \leq k) &= \mathbb P(X + Y \leq k \cap (X > Y \cup X \leq Y)) \\
    &= \mathbb P(X + Y \leq k)<
    \end{align*}
    for all $k,$ implying $\min(X, Y) + \max(X, Y) = X + Y.$

    \item Write $X, Y \sim$ Bernoulli(1/2) be independent, with $M = \max(X, Y)$ and $m = \min(X, Y),$ yielding $\Cov{M}{m} = 1/16 \neq 0 = \Cov{X}{Y}.$ Thus, this is a counterexample. 

    \end{enumerate}

    \item \begin{enumerate}
        \item We have $\mathbb E[X]$ satisfying \begin{align*}
            \mathbb E[X] = \mathbb E[S|U|] &= 1/2  \cdot \mathbb E[S|U| \mid S = 1] + 1/2 \cdot \mathbb E[S|U| \mid S = -1] \\
            &=\dfrac{1}{2} \cdot \mathbb E[|U| - |U|] = 0.
        \end{align*}

        Similarly, we have $$\mathbb E[X^2] = \mathbb E[S^2 U^2] = \mathbb E[U^2] = \Var{U} + \mathbb E[U]^2 = \Var{U} = 1,$$
        meaning $\Var{X} = 1$ as required.

        \item We have $$
        \mathbb P(0 < X < x) = \dfrac{1}{2} \mathbb P(0 < U <x) + \dfrac{1}{2} \mathbb P(-x < U< 0),
        $$
        which by symmetry of the Guassian distribution, we have
        $$
        \mathbb P(0 < X< x) = 2 \cdot \dfrac{1}{2} \cdot \mathbb P(0 < U < x) = \mathbb P(0 < \mathcal N(0, 1)  < x).
        $$

        \item Computing $\mathbb E[XY],$ we have
        $$
        \mathbb E[XY] = \mathbb E[S^2 |U|||V] = \mathbb E[|U||V|] = \mathbb E[|U|] \mathbb [|V|] = \dfrac{2}{\pi}
        $$
        since $U, V$ independent, yielding
        $$
            \mathbb E[XY] = 2/\pi.
        $$

        Hence, we have $\Cov{X}{Y} = 2/\pi - (0)(0) = 2/\pi$ as required.

        \item The pair $(X, Y)$ have support on $[0, \infty)^2 \cup (-\infty, 0]^2 \neq \mathbb R^2,$ meaning they cannot be jointly normal, since the joint normal distribution has support on $\mathbb R^2.$
    \end{enumerate}
\end{enumerate}
\end{document}